\documentclass{article}
\begin{document}
\section*{Binary native coefficient equality from independently grouped flat charts}
The additive normalization theorem treats typed additive circuits built
from composition, the additive bifunctor, the unit comparisons at zero,
and the middle-four. It interprets addition through mathlib's symmetric
monoidal category API. The flat route skips formal zeros, preserves
leaves with genuine zero coefficients, concatenates blocks, and implements
each block exchange by successive adjacent crossings. The two chart maps
are independently grouped from their respective trees; neither is defined
using the other chart and the asserted commuting square.

At one fixed binary final target, the separate route distributes the
flat occurrence fibre, applies the tensor of the two complete dependent
cells, and shuffles at the target. The joint route uses the canonical
four-level interchange, groups on its own chart, shuffles at the leaves,
uses the actual independent continuation cells, and then uses the actual
independent parent cell. The Lean proof derives the equality of these
two coefficient composites by distributivity naturality and the previously
verified ragged equality. Its raw input word lists the addresses
$(x,u,y,v)$ with $u$ chosen from the branch at $x$ and $v$ from the branch
at $y$. The chosen joint order is $(x,y,u,v)$. Mapping the raw lists to
coefficient words preserves order and multiplicity, even when coefficient
labels agree; no component of any whole-fibre arrow is selected by equality
of labels. Empty axes and branches use the declared absorption, not
initiality. There is no additional coproduct hypothesis.

The rig input remains explicit. The standard additive part is a symmetric
monoidal structure, and multiplication has the declared bifunctor,
distributor and absorber data. The binary-grid assumptions are obtained
from the actual GridData equations with the adapter. Multiplicative
middle-four compatibility can be derived by the verified structural
linearity adapter from the stated primitive associator/braiding and
binary/nullary distributor laws. The final theorem takes these
resulting laws and naturalities as premises. It does not package
a single library definition of a complete Laplaza rig.

The flat permutation program is computed from the chosen typed structural
circuit. Equality with every independently written program that has the
same permutation is not asserted here. Nor is identification with an
unrelated preinstalled enumeration of an abstract finite set asserted
without the address reindexing. This is fixed-final-fibre coefficient
naturality at arity two. Total augmentation transport, global variation of
the native maps, dependent contextad coherence and all finite occurrence
arities remain separate obligations.
\end{document}
